\documentclass[11pt]{article}

\usepackage[a4paper,margin=1in]{geometry}
\usepackage{amsmath}
\usepackage{amssymb}
\usepackage{amsthm}
\usepackage{booktabs}
\usepackage{graphicx}
\usepackage{natbib}
\usepackage{microtype}
\usepackage[hidelinks]{hyperref}

% ---------------------------------------------------------------------------
% Placeholder macro for results not yet computed (visible in the compiled PDF).
% ---------------------------------------------------------------------------
\newcommand{\TODO}[1]{\textbf{[TODO: #1]}}

% ---------------------------------------------------------------------------
% Notation shortcuts.
% ---------------------------------------------------------------------------
\newcommand{\E}{\mathbb{E}}
\newcommand{\Prob}{\mathbb{P}}
\newcommand{\Qmeas}{\mathbb{Q}}
\newcommand{\R}{\mathbb{R}}
\newcommand{\ois}{\mathrm{ois}}
\newcommand{\key}{\mathrm{key}}
\newcommand{\dd}{\,\mathrm{d}}
\newcommand{\Fwd}{F}
\newcommand{\Nvol}{\sigma_{\mathrm N}}
\newcommand{\Nmkt}{\sigma^{\mathrm{mkt}}}
\newcommand{\Nmdl}{\sigma^{\mathrm{mdl}}}

\theoremstyle{plain}
\newtheorem{proposition}{Proposition}
\theoremstyle{definition}
\newtheorem{remark}{Remark}

\title{\bfseries A multi-curve affine stochastic-volatility model for
RFR caplet surfaces: calibration and hard no-arbitrage projection}

\author{Andrey Popov \and Roland Grinis\thanks{Draft manuscript prepared for
\emph{Risk} (Cutting Edge). Corresponding author: andreykitsune@gmail.com.
Market data sourced from Bloomberg (VCUB caplet volatilities and ICVS curves,
as of 30 June 2026); redistribution of the raw data is not permitted.}}

\date{This draft: \today}

\begin{document}
\maketitle

\begin{abstract}
\noindent
We develop and calibrate a multi-curve affine stochastic-volatility short-rate
model for the normal-implied-volatility surface of per-period RFR and IBOR
caplets. The starting point is a structural observation, stated as a
proposition and independent of any particular market: a single-factor Gaussian
short-rate model implies a strike-independent Bachelier smile, so reproducing a
non-flat smile \emph{requires} stochastic volatility. We meet this requirement
with a two-factor lifted-CIR Hull--White--Heston specification with an
Ornstein--Uhlenbeck basis factor, and we position it explicitly as a
finite-dimensional affine realization of the Heath--Jarrow--Morton multi-curve
framework. Calibration combines a composite normal-volatility objective, a
zeroth-order Adam--SPSA optimiser robust to Monte Carlo noise, GPU-vectorised
simulation, and a hard projection onto the static no-arbitrage cone by pool
adjacent violators. We report the calibration protocol and an empirical
programme on EUR, GBP and USD surfaces as of 30~June~2026, including a
nested-model hierarchy that isolates the marginal contribution of each
modelling ingredient.
\end{abstract}

% ===========================================================================
\section{Motivation}
% ===========================================================================

The transition from IBOR-style term rates to overnight risk-free rates (RFRs),
together with the entrenchment of multi-curve discounting, has reshaped the
interest-rate modelling problem. Three structural facts now organise any
serious specification. First, discounting and forwarding are distinct: a single
yield curve can no longer serve both, and OIS (collateral) discounting must be
separated from the projection of the payoff-relevant index
\citep{Henrard2014,GrbacRunggaldier2016}. Second, RFR benchmarks are
backward-looking averages of an overnight rate, so caplet payoffs depend on a
compounded or averaged rate that is only known at the end of the accrual period,
and central-bank meetings inject scheduled, jump-like behaviour into the
policy-linked component \citep{FGGS2020,FGS2023}. Third, the option-implied
volatility surfaces of these markets display pronounced smile and skew that a
classical one-factor Gaussian short-rate model cannot reproduce.

This last point is not a calibration inconvenience but a structural
impossibility, and it is the analytical hook of the paper. If the
payoff-relevant rate is an affine function of a Gaussian factor---as it is in
any single-factor Hull--White model, however elaborate its deterministic
volatility or mean-reversion term structure---then the caplet is priced by the
Bachelier formula with a \emph{single} normal volatility that does not depend on
the strike. The model-implied smile is therefore flat by construction. We make
this precise as Proposition~\ref{prop:flat} below. Its practical corollary is a
warning against a common modelling error: one cannot manufacture strike
curvature by over-parametrising a Gaussian model in time. A smile is evidence of
stochastic volatility, and it must be modelled as such.

We take three deliberate steps in response, and they order the paper. (i) We
adopt a two-factor lifted-CIR Hull--White--Heston (HWH) specification with an
Ornstein--Uhlenbeck (OU) basis factor for the OIS/key-rate spread; the Heston
variance is the minimal economically meaningful ingredient that generates the
smile, and the ``lifting'' of the CIR variance into two Markovian factors
\citep{AbiJaber2019} adds term-structure flexibility at negligible cost. (ii) We
position the finite-dimensional short-rate construction inside the
Heath--Jarrow--Morton (HJM) multi-curve framework of
\citet{CFG2016,CFG2019}, stating which forward-rate dynamics the short-rate
factors imply and which no-arbitrage drift restrictions are inherited. (iii) We
treat static no-arbitrage of the fitted surface as a first-class constraint,
combining soft penalties during calibration with a hard projection onto the
butterfly/calendar cone by the pool-adjacent-violators algorithm (PAVA).

The computational machinery---GPU-vectorised Monte Carlo pricing and a
zeroth-order Adam--SPSA optimiser---is what makes the chosen model calibratable
against a noisy, non-smooth objective. We describe it in full, but we do not
present it as a separate theoretical contribution: its role is to enable the
model, not to be the model \citep{Savine2020,CapriottiGiles2024}. The empirical
programme is developed on EUR (a genuine dual-curve case, ESTR discounting with
6M EURIBOR projection), with GBP (SONIA, single-curve) and USD (SOFR discounting
with a Fed Funds ``policy'' curve and a market-quoted basis) as robustness
checks, all as of 30~June~2026.

% ===========================================================================
\section{The multi-curve affine model and its HJM positioning}
% ===========================================================================

\subsection{State dynamics}

The model carries a four-dimensional Markov state $(x_t, k_t, V^1_t, V^2_t)$
under the risk-neutral measure $\Qmeas$. Two Heston/CIR variance factors drive
a common stochastic volatility; a Hull--White factor $x$ drives the level of the
risk-free short rate; an OU factor $k$ drives the basis between the risk-free
(OIS) curve and the key/policy curve:
\begin{align}
\dd V^i_t &= \kappa_i\bigl(\theta_i(t) - V^i_t\bigr)\dd t
             + \varepsilon\,\sqrt{V^i_t}\,\dd W^{v_i}_t, \qquad i = 1,2,
             \label{eq:cir}\\
V_t &= w\,V^1_t + (1-w)\,V^2_t, \label{eq:varmix}\\
\dd x_t &= -\lambda\,x_t\,\dd t + \sqrt{V_t}\,\dd W^{x}_t, \label{eq:hw}\\
\dd k_t &= -\gamma\,k_t\,\dd t + \xi\,\dd W^{k}_t, \label{eq:ou}
\end{align}
with instantaneous leverage $\dd\langle W^{x}, W^{v_i}\rangle_t = \rho(t)\dd t$
between the rate and variance shocks; $W^{k}$ is taken independent of the rate
and variance drivers. The two forward curves are recovered by deterministic
shifts that fit the initial term structures exactly:
\begin{align}
r^{\ois}_t &= f^{\ois}(t) + x_t, \label{eq:rois}\\
r^{\key}_t &= f^{\key}(t) + x_t + s_{\mathrm{adj}}(t) + k_t. \label{eq:rkey}
\end{align}
Here $f^{\ois}(\cdot)$ and $f^{\key}(\cdot)$ are the initial instantaneous
forward curves for the OIS and key/policy indices, $s_{\mathrm{adj}}(\cdot)$ is a
deterministic spread adjustment, and $k_t$ is the stochastic basis. Discounting
is always done off the OIS curve; the key rate \eqref{eq:rkey} shares the
risk-free level $x_t$ but carries its own forward $f^{\key}$, deterministic
adjustment and stochastic basis $k_t$. For USD, $k_t$ is the Fed~Funds--SOFR
policy basis; for EUR it plays the role of the ESTR--EURIBOR tenor basis; for
single-curve GBP (SONIA) it is switched off, $\xi = 0$.

The mixture \eqref{eq:varmix} is a two-node Markovian ``lift'' of a CIR variance
in the sense of \citet{AbiJaber2019}: a sum of ordinary CIR factors with
distinct mean-reversion speeds $\kappa_1 \neq \kappa_2$ approximates a
richer---in the limit, fractional---variance kernel while keeping the state
finite-dimensional and affine. Setting $\varepsilon = 0$ collapses
\eqref{eq:cir} to a deterministic variance and \eqref{eq:hw} to a Gaussian
Hull--White factor: the vol-of-vol $\varepsilon$ is precisely the Heston
ingredient whose absence forces a flat smile (Proposition~\ref{prop:flat}).

The free parameter vector $\Theta$ has 27 components: eight PCHIP spline nodes
for each of the two variance targets $\theta_1(\cdot), \theta_2(\cdot)$ (16
nodes), together with 11 scalar and low-dimensional parameters---the mean
reversions $\lambda, \gamma$, the variance speeds $\kappa_1, \kappa_2$, the
vol-of-vol $\varepsilon$, the basis volatility $\xi$, the mixing weight $w$, and
the parameters of the leverage function $\rho(\cdot)$ and spread
$s_{\mathrm{adj}}(\cdot)$. The targets $\theta_i(t)$ and the correlation
$\rho(t)$ are represented by shape-preserving PCHIP splines, which keeps the
implied volatility term structure smooth without introducing spurious
oscillation.

\subsection{The Gaussian impossibility result}

The following statement is market-independent: it survives the move from the
Russian key-rate market of the original thesis to the developed-market surfaces
studied here, because it is a statement about the Gaussian model class, not about
any dataset.

\begin{proposition}[Flat Bachelier smile under Gaussian dynamics]
\label{prop:flat}
Fix a caplet on a period rate $A$ over $[T_{\mathrm{fix}},
T_{\mathrm{fix}}+\tau]$ settled at $T_{\mathrm{pay}}$, with time-$0$ value
$V_0(T,K) = D(T_{\mathrm{pay}})\,\tau\,\E^{\Qmeas^{T_{\mathrm{pay}}}}
[(A-K)^+]$ under the $T_{\mathrm{pay}}$-forward measure. Suppose $A$ is, under
$\Qmeas^{T_{\mathrm{pay}}}$, an affine function of a Gaussian process---as it is
whenever $\varepsilon = 0$ in \eqref{eq:cir}--\eqref{eq:rkey}, i.e.\ in any
single-factor Gaussian Hull--White model with deterministic (possibly
time-dependent) volatility and mean reversion. Then $A \sim
\mathcal{N}(\Fwd, s^2)$ with $\Fwd = \E^{\Qmeas^{T_{\mathrm{pay}}}}[A]$ and a
strike-independent variance $s^2$, and
\begin{equation}
\E^{\Qmeas^{T_{\mathrm{pay}}}}[(A-K)^+]
   = s\,\phi\!\Bigl(\tfrac{\Fwd-K}{s}\Bigr)
   + (\Fwd-K)\,\Phi\!\Bigl(\tfrac{\Fwd-K}{s}\Bigr),
\label{eq:bachelier}
\end{equation}
which is the Bachelier price with normal volatility $\Nvol = s/\sqrt{t}$ for the
relevant expiry $t$. The implied normal volatility $K \mapsto \Nvol(K)$ is
therefore constant in $K$: the model-implied smile is flat. Consequently no
choice of deterministic volatility or mean-reversion term structure inside the
Gaussian class reproduces a non-flat normal-volatility smile; strike curvature
requires the conditional variance of $A$ to be stochastic, i.e.\ $\varepsilon
> 0$.
\end{proposition}

\begin{proof}[Proof sketch]
Under $\Qmeas^{T_{\mathrm{pay}}}$ the state $x$ in \eqref{eq:hw} with
$\varepsilon=0$ is a Gaussian Ornstein--Uhlenbeck process, so any affine
functional $A$ of its path is Gaussian with a variance $s^2$ determined by the
(deterministic) integrated volatility and the payoff averaging, and in
particular independent of the strike $K$. Evaluating $\E[(A-K)^+]$ for
$A\sim\mathcal N(\Fwd,s^2)$ gives \eqref{eq:bachelier}. Matching to the Bachelier
functional form identifies a unique $\Nvol$ through $s = \Nvol\sqrt{t}$, the same
$s$ for every $K$; hence $\Nvol(K)$ is constant. The final claim follows because
strike dependence of $\Nvol$ requires the conditional law of $A$ to be
non-Gaussian, which \eqref{eq:cir}--\eqref{eq:hw} deliver only for
$\varepsilon>0$ through the mixing over the stochastic variance path.
\end{proof}

Proposition~\ref{prop:flat} motivates stochastic volatility as a
\emph{necessary} ingredient rather than an optional refinement, and it justifies
the minimal-extension philosophy of the model: we add exactly the degrees of
freedom---a stochastic, leverage-correlated variance---needed to break the
Gaussian degeneracy, and no more.

\subsection{Positioning in the HJM multi-curve framework}

A finite-dimensional short-rate model is best understood not as a computational
shortcut but as a tractable realization of a no-arbitrage HJM term-structure
model \citep{ChiarellaKwon2003}. We make the reduction explicit.

\paragraph{Implied OIS forward dynamics.}
Because $(x, V^1, V^2)$ is an affine process, the OIS bond price is
exponential-affine in the state,
\begin{equation}
P^{\ois}(t,T) = \exp\!\Bigl(-\!\int_t^T\! f^{\ois}(u)\dd u
   \;-\; B(t,T)\,x_t \;-\; \Gamma(t,T;V_t)\Bigr),
   \qquad
   B(t,T) = \frac{1-e^{-\lambda(T-t)}}{\lambda},
\label{eq:bond}
\end{equation}
where $\Gamma$ solves the variance Riccati system associated with the affine
generator of \eqref{eq:cir}--\eqref{eq:hw} and $\Gamma(t,T;0)$ recovers the
deterministic-volatility Hull--White convexity term. Differentiating
$f^{\ois}(t,T) = -\partial_T \log P^{\ois}(t,T)$ shows that the model is an HJM
model with a stochastic, separable instantaneous-forward volatility
\begin{equation}
\Sigma^{\ois}(t,T) = \sqrt{V_t}\;e^{-\lambda(T-t)},
\label{eq:hjmvol}
\end{equation}
i.e.\ the classical Hull--White exponential factor loading scaled by the
stochastic level $\sqrt{V_t}$. The HJM no-arbitrage drift restriction
\begin{equation}
\alpha^{\ois}(t,T) = \Sigma^{\ois}(t,T)\int_t^T \Sigma^{\ois}(t,u)\dd u
\label{eq:hjmdrift}
\end{equation}
is inherited automatically: it is exactly the condition encoded in the affine
convexity term $\Gamma$, and fitting $f^{\ois}$ to the initial curve enforces
the $t=0$ consistency $P^{\ois}(0,T) = \exp(-\int_0^T f^{\ois}(u)\dd u)$. We
therefore \emph{impose} no drift by hand; the drift restriction is a consequence
of specifying the short rate as \eqref{eq:rois} with an affine variance.

\paragraph{Multi-curve spread specification.}
The key/policy curve \eqref{eq:rkey} is an \emph{additive short-rate spread}
over OIS, driven by the deterministic $s_{\mathrm{adj}}$ and the stochastic OU
factor $k$. In the taxonomy of affine multiple-yield-curve models of
\citet{CFG2019}, this places the model in the ``exogenous spread'' class: the
risk-free curve and the spread are jointly affine, the spread is
mean-reverting and---because $\xi>0$ and $W^k$ is a genuine driver---stochastic,
and the tenor/key bond $P^{\key}(t,T)$ is exponential-affine with an extra
loading on $k_t$. This is the finite-dimensional counterpart of the
infinite-dimensional multi-curve HJM constructions of
\citet{MoreniPallavicini2014, PallaviciniTarenghi2010, SabelliEtal2018,
BienekEtal2023, MantiMolteni2023}; it retains the economically essential
separation of discounting and forwarding while remaining calibratable on the
sparse caplet data actually available. Compared with L\'evy-HJM
\citep{CrepeyEtal2015} or Wishart multi-curve specifications, the present model
is deliberately more modest and more transparent: a two-factor lifted-CIR
variance is rich enough to fit the observed smile and simple enough to diagnose.

% ===========================================================================
\section{Instrument and pricing}
% ===========================================================================

\subsection{Per-period caplets}

The calibration instrument is the per-period caplet. For an RFR (backward-looking)
index, the period rate $A$ over $[T_{\mathrm{fix}}, T_{\mathrm{fix}}+\tau]$ is the
compounded/averaged overnight rate; for the 6M~EURIBOR tenor it is the
forward-looking term rate; in both cases the caplet pays $\tau\,(A-K)^+$ at
$T_{\mathrm{pay}} = T_{\mathrm{fix}}+\tau$. Its time-zero value under the
$T_{\mathrm{pay}}$-forward measure is
\begin{equation}
V_0(T,K) = D(T_{\mathrm{pay}})\;\tau\;
   \E^{\Qmeas^{T_{\mathrm{pay}}}}\!\bigl[(A-K)^+\bigr],
\label{eq:caplet}
\end{equation}
with $D(\cdot)$ the OIS discount factor. We price \eqref{eq:caplet} by Monte
Carlo, simulating the joint state \eqref{eq:cir}--\eqref{eq:ou} and forming the
period rate consistently with the index convention.

\begin{remark}[Instrument/model consistency]
This is the developed-market reformulation of the calibration target. The
earlier Russian-market study approximated the payoff by a whole-period
\emph{average} rate; here the instrument is the genuine per-period caplet, and
the Monte Carlo pricer simulates the period rate directly. Instrument and model
are thus aligned, which removes a systematic pricing bias at short maturities
and---as we discuss in Section~\ref{sec:results}---changes the correct expiry
scaling of the Bachelier inversion.
\end{remark}

\subsection{Normal-volatility inversion}

Interest-rate caplets in the current regime are quoted in Bachelier (normal)
implied volatility, which we adopt as the calibration currency. Writing the
undiscounted Bachelier price as
\begin{equation}
\mathcal B(\Fwd, K, \nu) = \nu\,\phi\!\Bigl(\tfrac{\Fwd-K}{\nu}\Bigr)
   + (\Fwd-K)\,\Phi\!\Bigl(\tfrac{\Fwd-K}{\nu}\Bigr),
   \qquad \nu = \Nvol\,\sqrt{t_{\mathrm{fix}}},
\label{eq:bach_undisc}
\end{equation}
the model normal volatility is the $\Nvol$ solving $D(T_{\mathrm{pay}})\,\tau\,
\mathcal B(\Fwd,K,\Nvol\sqrt{t_{\mathrm{fix}}}) = V_0(T,K)$. The relevant Bachelier
expiry is the \emph{fixing time} $t_{\mathrm{fix}}$ of the caplet, giving the
$\sqrt{t_{\mathrm{fix}}}$ scaling in \eqref{eq:bach_undisc}; this replaces the
$\sqrt{T/3}$ average-rate scaling appropriate to a whole-period-average payoff
and is the correct convention once instrument and model agree. We invert by
Newton--Raphson on $\Nvol$, using the closed-form vega $\partial_\nu \mathcal B =
\phi((\Fwd-K)/\nu)$ and hence $\partial_{\Nvol}(D\tau\mathcal B) = D\tau
\sqrt{t_{\mathrm{fix}}}\,\phi(d)$, with a bracketed Brent fallback for the deep
wings where vega is small. Prices at or below intrinsic invert to zero
volatility and are excluded from the fit.

% ===========================================================================
\section{Calibration and no-arbitrage}
\label{sec:calib}
% ===========================================================================

\subsection{Composite objective}

We calibrate directly on the normal-volatility surface. Let $i$ index the caplet
grid nodes $(K_i, T_i)$ with market normal vols $\Nmkt_i$ and
model normal vols $\Nmdl_i(\Theta)$. The core loss is a
maturity-bucketed, strike-weighted root-mean-square vol error,
\begin{equation}
\mathcal L_{\mathrm{fit}}(\Theta)
   = \sum_{j}\, m_j \sqrt{
       \frac{\sum_{i\in \mathcal B_j} s_i\,
            \bigl(\Nmdl_i(\Theta) - \Nmkt_i\bigr)^2}
            {\sum_{i\in \mathcal B_j} s_i}},
\label{eq:coreloss}
\end{equation}
where $\mathcal B_j$ is the set of nodes in maturity bucket $j$, $m_j$ are
maturity weights, and $s_i$ are strike weights that de-emphasise deep-ITM nodes
(a smaller weight for deep-in-the-money strikes, unit weight for ATM/OTM),
renormalised within each bucket. Bucketing by maturity prevents the long-dated,
high-vol nodes from dominating a plain global RMSE and gives every expiry a
comparable voice in the fit. The full objective adds a quadratic
roughness penalty on the $\theta$-node values, an optional discrete $H^1$
(Sobolev) semi-norm on $\log\Nvol(K,T)$ that discourages surface roughness
everywhere \citep{Gatheral2006}, and the soft no-arbitrage penalties of
Section~\ref{sec:noarb}:
\begin{equation}
\mathcal L(\Theta) = \mathcal L_{\mathrm{fit}}(\Theta)
   + \lambda_\theta \sum_n (\Delta\theta_n)^2
   + \lambda_S\, \mathcal L_{\mathrm{Sob}}(\Theta)
   + \lambda_{\mathrm{arb}}\, \mathcal P_{\mathrm{arb}}(\Theta).
\label{eq:totalloss}
\end{equation}

\subsection{Adam--SPSA under Monte Carlo noise}

The objective \eqref{eq:totalloss} is (i) noisy, because
$\Nmdl_i$ is a Monte Carlo estimate, and (ii) non-smooth, because
the arbitrage penalty and the vol inversion introduce kinks. A pathwise or
adjoint gradient of the full pipeline is therefore both expensive and unreliable
near those kinks---the estimated gradient tends to collapse or blow up. We
instead use simultaneous perturbation stochastic approximation (SPSA), a
black-box zeroth-order estimator that requires only two objective evaluations per
step regardless of dimension. With a Rademacher perturbation vector $\Delta$
(components $\pm 1$) and step $c$, the gradient estimate is
\begin{equation}
\widehat g_n(\Theta) = \frac{\mathcal L(\Theta + c\Delta)
   - \mathcal L(\Theta - c\Delta)}{2c}\;\frac{1}{\Delta_n},
   \qquad n = 1,\dots,27,
\label{eq:spsa}
\end{equation}
which we feed into an Adam update. The common-random-numbers device (the same
Monte Carlo seed for the $\Theta\pm c\Delta$ evaluations) suppresses the noise in
the finite difference. This scheme is robust to both the MC noise and the PAVA
kinks and avoids the gradient collapse of a naive differentiate-through-inversion
approach.

\subsection{GPU vectorisation as an enabler}

Monte Carlo calibration of an affine stochastic-volatility model over a full
$(K,T)$ grid, with repeated no-arbitrage diagnostics inside the optimiser loop,
is only feasible at interactive speed if pricing is vectorised. We simulate
paths as the batch dimension on the GPU and accumulate the option payoffs with
running sums, so the memory footprint is $O(\text{paths})$ rather than
$O(\text{paths}\times\text{steps})$ and a full surface reprice is a handful of
tensor operations. This is an implementation enabler, not a separate theoretical
contribution: its value is that it makes the chosen model calibratable and the
nested-model programme of Section~\ref{sec:results} practical
\citep{Savine2020,CapriottiGiles2024}. Adjoint (AAD) sensitivities for caplet
Greeks and calibration Jacobians are a natural future refinement but are not
required here.

\subsection{Hard no-arbitrage projection by PAVA}
\label{sec:noarb}

Soft penalties reduce but do not eliminate static-arbitrage violations. We
therefore add a hard projection step: the fitted surface is projected onto the
no-arbitrage cone by alternating two elementary convex projections. This is the
part of the pipeline we regard as most important for producing a usable surface.

\paragraph{Butterfly (per maturity slice).}
For each maturity $T$ we convert the slice of normal vols to undiscounted
Bachelier prices $C(K)$ and enforce the strike-convexity conditions that
characterise butterfly arbitrage-freeness: $C(K)$ non-increasing and convex in
$K$, and bounded below by intrinsic. On the non-uniform strike grid the discrete
convexity constraint reads
\begin{equation}
C(K_i) \;\le\; \frac{\Delta K_i\, C(K_{i-1}) + \Delta K_{i-1}\, C(K_{i+1})}
                    {\Delta K_{i-1} + \Delta K_i},
\qquad \Delta K_i = K_{i+1} - K_i,
\label{eq:convex}
\end{equation}
alongside $C(K_i)\ge C(K_{i+1})$ and $C(K_i)\ge (\Fwd-K_i)^+ D\tau$. We enforce
\eqref{eq:convex} and the monotonicity by an $O(n)$-per-pass alternating scan
(a few passes suffice), then invert the repaired prices back to normal vols by
the Newton step of \eqref{eq:bach_undisc}. Deep-OTM strikes whose repaired price
inverts below a floor are left at their current value rather than zeroed.

\paragraph{Calendar (per strike).}
For each strike $K$ we require the total variance $w(T) = \Nvol(K,T)^2\,T$ to be
non-decreasing in maturity---the calendar no-arbitrage condition. We project the
sequence $\{w(T)\}$ onto the non-decreasing cone by isotonic regression (PAVA),
which returns the $L^2$-nearest non-decreasing sequence, and map back to vols by
$\Nvol = \sqrt{w/T}$.

Both maps are Euclidean projections onto closed convex sets, so alternating them
(butterfly across strikes within a slice, calendar across maturities within a
strike) converges to a point in their intersection by the alternating-projection
/ POCS argument; a short final cleanup loop guarantees both constraints hold
simultaneously at every grid node. The projection is minimal in the sense that
it perturbs the fitted surface by the smallest $L^2$ amount consistent with
static no-arbitrage.

\subsection{Decomposing the source of arbitrage}
\label{sec:arbsrc}

A projection tells us \emph{where} a surface violates no-arbitrage but not
\emph{why}. Following the referee guidance, we separate three sources, because
each calls for a different remedy. (a) \emph{Market/stripping data}: the flat
cap-vol quotes must be stripped into per-period caplet vols, and quote
granularity or bid--offer noise can inject small non-convexities in the
resulting surface---remedied by pre-smoothing the input quotes. (b) \emph{Monte
Carlo noise}: finite-path pricing perturbs each node independently, which can
create spurious local violations---remedied by more paths, antithetic variates or
control variates, and diagnosable because it shrinks with the path count. (c)
\emph{Parametric model}: the affine family may be unable to hit every node
simultaneously, so a residual, \emph{systematic} violation survives even in the
noise-free limit---this is the genuine model-limitation term, and it is what the
PAVA projection legitimately repairs. Section~\ref{sec:results} reports the
violation count attributable to each source by re-pricing at increasing path
counts and comparing against a pre-smoothed input surface.

% ===========================================================================
\section{Empirical results}
\label{sec:results}
% ===========================================================================

\emph{This section is a structured skeleton. The market-agnostic protocol is
fixed; the numerical results of the calibration campaign are pending and are
marked with \TODO{} placeholders.}

\subsection{Data and caplet stripping}

We use Bloomberg VCUB flat cap volatilities and ICVS discount/forward curves as
of 30~June~2026 for EUR, GBP and USD. EUR is the primary, genuinely dual-curve
case: ESTR discounting with 6M~EURIBOR projection. GBP is single-curve SONIA.
USD combines SOFR discounting with a Fed~Funds ``policy'' curve and a
market-quoted Fed~Funds--SOFR basis, which the OU factor $k_t$ represents. Flat
cap volatilities are stripped into a per-period caplet surface by a bootstrap
that is exact on the input: the flat-cap reprice error of the stripped surface is
$0.000$~bp, so the stripping step introduces no bias into the calibration target.
The EUR ATM caplet normal volatilities run from roughly $48$~bp at the $6$M
expiry to $\sim 73$~bp near $5$Y and back to $\sim 66$~bp at $20$Y, with a convex
smile whose minimum sits near the forward---exactly the strike curvature that
Proposition~\ref{prop:flat} rules out for any Gaussian model.

\begin{table}[t]
\centering
\caption{EUR caplet surface: dimensions and ATM level by maturity bucket.}
\label{tab:data}
\begin{tabular}{lccc}
\toprule
Maturity bucket & \# strikes & ATM normal vol (bp) & Forward (\%) \\
\midrule
$\le 1$Y   & \TODO{n} & \TODO{ATM} & \TODO{fwd} \\
$1$--$5$Y  & \TODO{n} & \TODO{ATM} & \TODO{fwd} \\
$5$--$10$Y & \TODO{n} & \TODO{ATM} & \TODO{fwd} \\
$>10$Y     & \TODO{n} & \TODO{ATM} & \TODO{fwd} \\
\bottomrule
\end{tabular}
\end{table}

\subsection{EUR calibration fit}

Figure~\ref{fig:surface} overlays the calibrated model surface on the market
surface. \TODO{EUR market-vs-model surface overlay figure}. We report the fit
with both absolute and relative metrics, following the referee's request that
publication-grade calibration not rely on absolute bp alone. Define the relative
error $e_i = (\Nmdl_i - \Nmkt_i)/\Nmkt_i$ and report, per maturity bucket, the
absolute RMSE (bp),
the relative RMSE (\%), and an ATM-versus-OTM decomposition. Table~\ref{tab:fit}
gives the shell.

\begin{table}[t]
\centering
\caption{EUR fit quality: absolute, relative, and ATM/OTM decomposition.}
\label{tab:fit}
\begin{tabular}{lcccc}
\toprule
Maturity bucket & Abs.\ RMSE (bp) & Rel.\ RMSE (\%) & ATM RMSE (bp) & OTM RMSE (bp) \\
\midrule
$\le 1$Y   & \TODO{} & \TODO{} & \TODO{} & \TODO{} \\
$1$--$5$Y  & \TODO{} & \TODO{} & \TODO{} & \TODO{} \\
$5$--$10$Y & \TODO{} & \TODO{} & \TODO{} & \TODO{} \\
$>10$Y     & \TODO{} & \TODO{} & \TODO{} & \TODO{} \\
\midrule
All        & \TODO{EUR overall RMSE} & \TODO{} & \TODO{} & \TODO{} \\
\bottomrule
\end{tabular}
\end{table}

Figure~\ref{fig:termstruct} shows the calibrated variance targets
$\theta_1(t),\theta_2(t)$ and leverage $\rho(t)$. \TODO{$\theta(t)$, $\rho(t)$
term-structure figure}.

\subsection{Nested-model hierarchy}
\label{sec:nested}

To isolate the marginal contribution of each modelling ingredient we calibrate a
nested hierarchy on the \emph{identical} EUR surface: (a) one-factor Gaussian
Hull--White; (b) two-factor Gaussian Hull--White; (c) Hull--White with a
deterministic volatility term structure; and (d) the proposed two-factor
lifted-CIR HWH. Models (a)--(c) all satisfy Proposition~\ref{prop:flat} (or a
term-structure-only relaxation of it) and cannot fit strike curvature; the ladder
quantifies exactly how much of the residual each ingredient removes and confirms
that the smile fit is attributable to stochastic volatility rather than to added
parameters per se. Table~\ref{tab:nested} is the results shell.

\begin{table}[t]
\centering
\caption{Nested-model ladder on the EUR surface (same data, same grid).}
\label{tab:nested}
\begin{tabular}{llccc}
\toprule
Model & \# params & Vol RMSE (bp) & Rel.\ RMSE (\%) & Smile captured? \\
\midrule
(a) 1F Gaussian HW               & \TODO{} & \TODO{} & \TODO{} & No (flat) \\
(b) 2F Gaussian HW               & \TODO{} & \TODO{} & \TODO{} & No (flat) \\
(c) HW + deterministic vol       & \TODO{} & \TODO{} & \TODO{} & Term-structure only \\
(d) 2F lifted-CIR HWH (proposed) & $27$    & \TODO{} & \TODO{} & Yes \\
\bottomrule
\end{tabular}
\end{table}

\subsection{No-arbitrage diagnostics and source decomposition}

We report static-arbitrage coverage before and after the PAVA projection of
Section~\ref{sec:noarb}, and the fit cost of the projection (the vol RMSE change
induced by the projection). Figure~\ref{fig:pava} shows a representative slice
before and after repair. \TODO{PAVA before/after figure}. Table~\ref{tab:arb}
decomposes the pre-projection violations into the three sources of
Section~\ref{sec:arbsrc}---market/stripping data, Monte Carlo noise, and
parametric model---by re-pricing at increasing path counts and against a
pre-smoothed input surface.

\begin{table}[t]
\centering
\caption{No-arbitrage coverage and decomposition of violation sources.}
\label{tab:arb}
\begin{tabular}{lcc}
\toprule
Quantity & Before PAVA & After PAVA \\
\midrule
Butterfly-arb-free nodes (\%)  & \TODO{} & \TODO{} \\
Calendar-arb-free nodes (\%)   & \TODO{} & \TODO{} \\
Vol RMSE (bp)                  & \TODO{} & \TODO{} \\
\midrule
\multicolumn{3}{l}{\emph{Violations attributable to:}} \\
\quad Market/stripping data    & \multicolumn{2}{c}{\TODO{count / \%}} \\
\quad Monte Carlo noise        & \multicolumn{2}{c}{\TODO{count / \%}} \\
\quad Parametric model         & \multicolumn{2}{c}{\TODO{count / \%}} \\
\bottomrule
\end{tabular}
\end{table}

\subsection{Short-maturity residual and a scheduled-jump study}
\label{sec:jump}

The continuous-diffusion model cannot fully represent scheduled central-bank
decisions, which is expected to leave a systematic residual at short maturities
($T < 6$M) where meeting risk dominates. We quantify this residual and test how
much of it is attributable to policy-event risk with a simple scheduled-jump
extension: at the known ECB (and, for GBP/USD, BoE/Fed) meeting dates in the
front window we augment \eqref{eq:rkey} with a compound-Poisson jump of small
fixed intensity concentrated on those dates,
\begin{equation}
\dd r^{\key}_t \;\longmapsto\; \dd r^{\key}_t + \dd J_t,
\qquad J_t = \sum_{m:\,t_m \le t} \eta_m,
\label{eq:jump}
\end{equation}
with $t_m$ the scheduled meeting dates and $\eta_m$ i.i.d.\ jump sizes. This is
an affine jump-diffusion in the sense of the stochastic-discontinuity literature
\citep{FGGS2020,FGS2023}; even uncalibrated, it lets us measure how much of the
$T<6$M residual smile the meeting schedule can plausibly explain.
Table~\ref{tab:jump} reports the short-end residual with and without the jump
term. \TODO{short-maturity residual + jump-study numbers}.

\begin{table}[t]
\centering
\caption{Short-maturity ($T<6$M) residual smile: diffusion vs.\ scheduled-jump.}
\label{tab:jump}
\begin{tabular}{lcc}
\toprule
Expiry & Diffusion residual RMSE (bp) & With scheduled jump (bp) \\
\midrule
$1$M & \TODO{} & \TODO{} \\
$3$M & \TODO{} & \TODO{} \\
$6$M & \TODO{} & \TODO{} \\
\bottomrule
\end{tabular}
\end{table}

\subsection{Cross-currency robustness and swaption out-of-sample}

We repeat the EUR protocol on GBP (SONIA) and USD (SOFR + Fed~Funds policy basis)
to check that the specification and calibration are not EUR-specific.
Table~\ref{tab:xccy} collects the headline fit metrics across currencies. As an
out-of-sample validation we price ATM swaptions---instruments not used in the
caplet calibration---from the calibrated caplet parameters and compare against
market ATM swaption vols; Figure~\ref{fig:swaption} shows the comparison.
\TODO{GBP/USD calibration numbers}; \TODO{swaption out-of-sample figure and
error}.

\begin{table}[t]
\centering
\caption{Cross-currency robustness: headline fit metrics.}
\label{tab:xccy}
\begin{tabular}{lcccc}
\toprule
Currency & Curve setup & Vol RMSE (bp) & Rel.\ RMSE (\%) & Arb-free after PAVA (\%) \\
\midrule
EUR & ESTR / 6M EURIBOR      & \TODO{} & \TODO{} & \TODO{} \\
GBP & SONIA (single-curve)   & \TODO{} & \TODO{} & \TODO{} \\
USD & SOFR / Fed Funds basis & \TODO{} & \TODO{} & \TODO{} \\
\bottomrule
\end{tabular}
\end{table}

\begin{figure}[t]
\centering
\framebox[0.85\textwidth]{\rule{0pt}{0.32\textwidth}\quad
\TODO{EUR market-vs-model caplet surface overlay}\quad}
\caption{Calibrated EUR caplet normal-volatility surface (model) overlaid on the
stripped market surface.}
\label{fig:surface}
\end{figure}

\begin{figure}[t]
\centering
\framebox[0.85\textwidth]{\rule{0pt}{0.28\textwidth}\quad
\TODO{$\theta_1(t),\theta_2(t),\rho(t)$ term structures}\quad}
\caption{Calibrated variance targets and leverage term structure.}
\label{fig:termstruct}
\end{figure}

\begin{figure}[t]
\centering
\framebox[0.85\textwidth]{\rule{0pt}{0.28\textwidth}\quad
\TODO{PAVA before/after slice}\quad}
\caption{A representative maturity slice before and after the hard PAVA
projection onto the butterfly/calendar cone.}
\label{fig:pava}
\end{figure}

\begin{figure}[t]
\centering
\framebox[0.85\textwidth]{\rule{0pt}{0.28\textwidth}\quad
\TODO{swaption ATM model-vs-market}\quad}
\caption{Out-of-sample ATM swaption volatilities: model (from calibrated caplet
parameters) vs.\ market.}
\label{fig:swaption}
\end{figure}

% ===========================================================================
\section{Conclusions and extensions}
% ===========================================================================

We have presented a multi-curve affine stochastic-volatility short-rate model
for RFR and IBOR caplet surfaces, organised around three ideas: a
market-independent proof that Gaussian short-rate models imply a flat Bachelier
smile and hence that stochastic volatility is necessary
(Proposition~\ref{prop:flat}); a two-factor lifted-CIR Hull--White--Heston
specification with a stochastic basis, positioned explicitly as a
finite-dimensional affine realization of the multi-curve HJM framework of
\citet{CFG2016,CFG2019}; and an arbitrage-aware calibration in which soft
penalties are followed by a hard PAVA projection onto the static no-arbitrage
cone. GPU-vectorised Monte Carlo and a zeroth-order Adam--SPSA optimiser are the
enablers that make the model calibratable against a noisy, non-smooth objective,
not contributions in their own right.

Several extensions follow naturally and are partly explored in
Section~\ref{sec:jump}. Scheduled affine jumps at policy-meeting dates are the
most economically motivated: the short-maturity residual is a symptom of the
continuous-diffusion assumption meeting discrete policy decisions, and the
stochastic-discontinuity framework \citep{FGGS2020,FGS2023} is the right next
step. Richer correlation and skew structure---a fuller leverage function or an
explicit event-risk skew component---would help distinguish diffusive from
event-driven asymmetry. Finally, pathwise and adjoint (AAD) sensitivities
\citep{CapriottiGiles2024,Savine2020} would turn the calibrated model into a
production risk engine for caplet and cap portfolios. None of these are
corrections of the approach; they are continuations of it.

% ===========================================================================
\section*{Acknowledgements}
% ===========================================================================
\noindent We thank Roland Grinis for detailed review comments that shaped the
structure of this paper. Market data: Bloomberg.

% ===========================================================================
\begin{thebibliography}{99}
% ===========================================================================

\bibitem[Abi Jaber(2019)]{AbiJaber2019}
Abi Jaber, E. (2019). Lifting the Heston model.
\emph{Quantitative Finance} 19(12), 1995--2013.

\bibitem[Bienek et al.(2023)]{BienekEtal2023}
Bienek, T., Deelstra, G., Lichtenstern, A. and Zagst, R. (2023).
A multi-curve HJM factor model for pricing and risk management.
\emph{Quantitative Finance} 23(10).

\bibitem[Cao et al.(2020)]{Cao2020}
Cao, J. et al. (2020). Calibration of stochastic-volatility models and the
implied-volatility surface. \TODO{confirm exact Cao calibration-benchmark reference}.

\bibitem[Capriotti and Giles(2024)]{CapriottiGiles2024}
Capriotti, L. and Giles, M. (2024). 15 years of Adjoint Algorithmic
Differentiation (AAD) in finance. \emph{Quantitative Finance}.

\bibitem[Chiarella and Kwon(2003)]{ChiarellaKwon2003}
Chiarella, C. and Kwon, O.K. (2003). Finite dimensional affine realisations of
HJM models in terms of forward rates and yields.
\emph{Review of Derivatives Research} 6, 129--155.

\bibitem[Cr\'epey et al.(2015)]{CrepeyEtal2015}
Cr\'epey, S., Grbac, Z., Ngor, N. and Skovmand, D. (2015). A L\'evy HJM
multiple-curve model with application to CVA computation.
\emph{Quantitative Finance} 15(3), 401--419.

\bibitem[Cuchiero et al.(2016)]{CFG2016}
Cuchiero, C., Fontana, C. and Gnoatto, A. (2016). A general HJM framework for
multiple yield curve modelling. \emph{Finance and Stochastics} 20(2), 267--320.

\bibitem[Cuchiero et al.(2019)]{CFG2019}
Cuchiero, C., Fontana, C. and Gnoatto, A. (2019). Affine multiple yield curve
models. \emph{Mathematical Finance} 29(2), 568--611.

\bibitem[Fontana et al.(2020)]{FGGS2020}
Fontana, C., Grbac, Z., Guembel, S. and Schmidt, T. (2020). Term structure
modelling for multiple curves with stochastic discontinuities.
\emph{Finance and Stochastics} 24(2), 465--511.

\bibitem[Fontana et al.(2023)]{FGS2023}
Fontana, C., Grbac, Z. and Schmidt, T. (2023). Term structure modeling with
overnight rates beyond stochastic continuity. \emph{Mathematical Finance}.

\bibitem[Gatheral(2006)]{Gatheral2006}
Gatheral, J. (2006). \emph{The Volatility Surface: A Practitioner's Guide}.
Wiley. (Total-variance / arbitrage-free surface conditions.)

\bibitem[Grbac and Runggaldier(2016)]{GrbacRunggaldier2016}
Grbac, Z. and Runggaldier, W. (2016). \emph{Interest Rate Modeling: Post-Crisis
Challenges and Approaches}. Springer.

\bibitem[Henrard(2014)]{Henrard2014}
Henrard, M. (2014). \emph{Interest Rate Modelling in the Multi-Curve Framework}.
Palgrave Macmillan.

\bibitem[Manti and Molteni(2023)]{MantiMolteni2023}
Manti, F. and Molteni, D. (2023). Parsimonious HJM-FMM model with the new
risk-free term rates. SSRN working paper.

\bibitem[Moreni and Pallavicini(2014)]{MoreniPallavicini2014}
Moreni, N. and Pallavicini, A. (2014). Parsimonious HJM modelling for multiple
yield curve dynamics. \emph{Quantitative Finance} 14(2), 199--210.

\bibitem[Pallavicini and Tarenghi(2010)]{PallaviciniTarenghi2010}
Pallavicini, A. and Tarenghi, M. (2010). Interest-rate modeling with multiple
yield curves. SSRN working paper.

\bibitem[Sabelli et al.(2018)]{SabelliEtal2018}
Sabelli, C., Pioppi, M., Sitzia, L. and Bormetti, G. (2018). Multi-curve HJM
modelling for risk management. \emph{Quantitative Finance} 18(4), 563--590.

\bibitem[Savine(2020)]{Savine2020}
Savine, A. (2020). \emph{Modern Computational Finance: AAD and Parallel
Simulations}. Wiley.

\end{thebibliography}

\end{document}
